\chapter{Diarrhoea}
\index{Diarrhoea}

\section*{Definition and Overview}
Diarrhoea is the passage of loose or watery stools more often than usual. It is a common symptom, usually caused by infection of the gut, but it can also result from foods, medicines, and other conditions. Most episodes are short-lived and settle on their own within a few days, and the main treatment is replacing lost fluids. Diarrhoea is most dangerous in babies, young children, older adults, and people with weakened immunity, because it can cause dehydration. This chapter covers the causes, symptoms, and management of diarrhoea.

\section*{Epidemiology}
Infectious diarrhoea is extremely common, affecting most people at some point each year, with young children affected most often. Gastroenteritis is a leading cause of GP visits and hospital presentations in children. Worldwide, diarrhoea remains a major cause of death in young children, mainly in low-income countries, although oral rehydration therapy and rotavirus vaccination have greatly reduced this toll. In many countries, outbreaks occur in childcare and aged care settings, and foodborne infection is a common cause.

\section*{Aetiology and Pathophysiology}
\begin{itemize}
    \item \textbf{Viral infection:} rotavirus, norovirus, and adenovirus are the most common causes
    \item \textbf{Bacterial infection:} salmonella, campylobacter, shigella, and E. coli from contaminated food or water
    \item \textbf{Parasites:} giardia and cryptosporidium, often from contaminated water
    \item \textbf{Food poisoning:} toxins in contaminated food cause rapid-onset diarrhoea and vomiting
    \item \textbf{Medicines:} antibiotics and some other drugs can cause diarrhoea
    \item \textbf{Chronic conditions:} irritable bowel syndrome, inflammatory bowel disease, coeliac disease, and food intolerances
    \item \textbf{Mechanism:} the gut loses its ability to absorb fluid, so water passes through in the stool
\end{itemize}

\section*{Clinical Features}
\begin{itemize}
    \item Loose, watery, or frequent stools
    \item Abdominal cramps, bloating, and urgency
    \item Nausea, vomiting, and fever in infectious cases
    \item Signs of dehydration: thirst, dry mouth, dark urine, dizziness, and fatigue
    \item Severe dehydration: confusion, sunken eyes, rapid heart rate, and reduced urine output
    \item Blood or mucus in the stool in some bacterial infections
    \item Most episodes settle within 2--5 days
\end{itemize}

\section*{Differential Diagnosis}
\begin{itemize}
    \item Viral versus bacterial gastroenteritis
    \item Food poisoning versus person-to-person infection
    \item Irritable bowel syndrome and functional diarrhoea
    \item Inflammatory bowel disease, with blood, mucus, and weight loss
    \item Coeliac disease and food intolerances
    \item Medication side effects
    \item Appendicitis and other abdominal emergencies in some presentations
\end{itemize}

\section*{When to Seek Medical Care}
See a doctor if diarrhoea lasts more than a few days, is severe, or is accompanied by high fever, blood in the stool, or signs of dehydration. Babies, young children, older adults, and people with chronic illness need earlier review.

\textbf{Contact your local emergency services for:}
\begin{itemize}
    \item Signs of severe dehydration: confusion, drowsiness, fainting, or no urine for many hours
    \item Blood in the stool with severe abdominal pain or high fever
    \item A baby who is listless, has a sunken fontanelle, or is not passing wet nappies
    \item Vomiting that prevents keeping down any fluids
\end{itemize}

\section*{Diagnosis and Investigations}
\begin{itemize}
    \item \textbf{History and examination:} most cases are diagnosed clinically
    \item \textbf{Stool tests:} for bacteria, parasites, or toxins when infection is severe, persistent, or part of an outbreak
    \item \textbf{Blood tests:} electrolytes and kidney function in significant dehydration
    \item \textbf{Assessment of dehydration:} clinical signs guide management
    \item \textbf{Chronic diarrhoea investigations:} tests for coeliac disease, inflammatory bowel disease, or intolerances
\end{itemize}

\section*{Management}
\begin{itemize}
    \item \textbf{Fluid replacement:} the mainstay of treatment; oral rehydration solution replaces water and salts
    \item \textbf{Continue eating:} bland foods as tolerated, including bread, rice, and bananas
    \item \textbf{Avoid aggravating foods:} rich, fatty, and very sugary foods
    \item \textbf{Medicines:} antidiarrhoeal agents are generally not recommended in infectious diarrhoea
    \item \textbf{Probiotics:} may shorten some episodes
    \item \textbf{Antibiotics:} only for specific bacterial infections
    \item \textbf{Hospital care:} intravenous fluids for severe dehydration
    \item \textbf{Hygiene:} hand washing to prevent spread to others
\end{itemize}

\section*{Complications}
\begin{itemize}
    \item Dehydration and electrolyte disturbance
    \item Acute kidney injury
    \item Malnutrition, especially in children
    \item Spread of infection within families and communities
    \item Post-infectious irritable bowel syndrome
    \item Rare complications such as haemolytic uraemic syndrome from E. coli infection
\end{itemize}

\section*{Prognosis and Outcomes}
Most episodes of diarrhoea resolve within a few days with fluids and rest. Complications are uncommon in healthy people when dehydration is prevented. With prompt treatment, babies, children, and older adults recover fully. Chronic causes of diarrhoea are managed by treating the underlying condition, and the outlook is generally good.

\section*{Prevention and Self-Care}
\begin{itemize}
    \item Wash hands thoroughly, especially after the toilet and before food
    \item Practise safe food handling and cooking
    \item Drink safe water, including when travelling
    \item Ensure rotavirus vaccination for infants
    \item Drink plenty of fluids during illness
    \item Stay home while infectious
    \item Seek early advice for vulnerable people
\end{itemize}

\section*{Sources and Further Reading}
The following reputable sources provide further information for healthcare students and patients seeking reliable, current advice about diarrhoea.

\begin{itemize}
    \item Healthdirect Australia. \textit{Diarrhoea: causes and treatment.} https://www.healthdirect.gov.au/
    \item NSW Health. \textit{Gastroenteritis and diarrhoea fact sheets.} https://www.health.nsw.gov.au/
    \item NHS. \textit{Diarrhoea: causes, treatment, and when to see a doctor.} https://www.nhs.uk/conditions/diarrhoea/
    \item Mayo Clinic. \textit{Diarrhoea: symptoms and treatment.} https://www.mayoclinic.org/
    \item MSD Manual. \textit{Diarrhoea: professional overview.} https://www.msdmanuals.com/
    \item World Health Organization. \textit{Diarrhoeal disease fact sheet.} https://www.who.int/
\end{itemize}
